% Work/MAT | 1. rocnik | nova hodina 32: Usmernovani zlomku.
% Dosavadni tabulka mocnin je uzavrena v hodine 30; nepotrebujeme novou.
% 31 castecnne odmocnovani; 32 jednoduche jmenovatele;
% 33 sdruzene vyrazy; ctvrtletni test z puvodni 32. hodiny se musi preplanovat.
% Ucebnice: Mocniny a Odmocniny.tex, oddil Usmernovani zlomku.
% Realisticky.cz 010209 Usmernovani zlomku: nasobeni jednickou a volba
% jednodussiho postupu. Ulohy vlastni, ne doslovny prepis.
\documentclass[11pt,a4paper]{article}
% Spolecny styl pro ucitelske pripravy MAT (Work/MAT)
% Verze 1.1, 2026-10-23
\NeedsTeXFormat{LaTeX2e}
\ProvidesPackage{priprava}[2026/10/23 v1.1 Sablona ucitelskych priprav MAT]

% Zakladni typografie a cestina
\RequirePackage[T1]{fontenc}
\RequirePackage[utf8]{inputenc}
\RequirePackage[czech]{babel}
\RequirePackage{lmodern}
\RequirePackage{microtype}
\RequirePackage{amsmath,amssymb}
\RequirePackage{xcolor}
\RequirePackage{enumitem}
\RequirePackage[most]{tcolorbox}
\RequirePackage{tabularx,booktabs,array}
\RequirePackage{fancyhdr}
\RequirePackage{needspace}
\RequirePackage{geometry}
\RequirePackage[hidelinks]{hyperref}

\geometry{a4paper,top=17mm,bottom=18mm,left=23mm,right=23mm}
\definecolor{notegray}{gray}{0.39}
\setlength{\parindent}{0pt}
\setlength{\parskip}{2.5pt plus 1pt minus 0.5pt}
\setlist[itemize]{topsep=2pt,itemsep=1pt,parsep=0pt,leftmargin=1.45em}
\setlist[enumerate]{topsep=2pt,itemsep=2pt,parsep=0pt,leftmargin=1.65em}
\setcounter{secnumdepth}{0}
\renewcommand{\arraystretch}{1.12}

% Zahlavi: vlevo rocnik | globalni tema, vpravo cislo: nazev hodiny
% Zapati zustava zcela prazdne.
\pagestyle{fancy}
\fancyhf{}
\renewcommand{\headrulewidth}{0.35pt}
\renewcommand{\footrulewidth}{0pt}
\setlength{\headheight}{14pt}
\newcommand{\nastavPripravu}[4]{%
  \fancyhead[L]{\small #1\enspace\textbar\enspace #2}%
  \fancyhead[R]{\small #3: #4}%
}

% Minimalisticke cernobile boxy. Kazda definice zvlast.
\tcbset{commonlesson/.style={enhanced,breakable=false,colback=white,colframe=black,
  colbacktitle=white,coltitle=black,boxrule=0.45pt,
  arc=0.8mm,outer arc=0.8mm,left=3mm,right=3mm,top=1.6mm,bottom=1.6mm,
  toptitle=1mm,bottomtitle=0.7mm,fonttitle=\bfseries,
  attach boxed title to top left={xshift=2.5mm,yshift=-1.15mm},
  boxed title style={colback=white,colframe=white,boxrule=0pt,left=1mm,right=1mm}}}
\newtcolorbox{zadaniBox}[1][]{commonlesson,title={Zadání},#1}
\newtcolorbox{prikladBox}[1][]{commonlesson,title={Příklad},#1}
\newtcolorbox{definiceBox}[1][]{commonlesson,title={Definice},#1}

% Spolecny box pro zadani a jeho reseni, bez mezery mezi castmi.
% Pouziti:
%   \begin{prikladReseniBox}
%     Zadani ulohy.
%     \tcblower
%     \begin{align*} ... \end{align*}
%   \end{prikladReseniBox}
% Spodni cast ma jemny sedy podklad a delici linku; lze v ni sazet
% vice radku matematiky. Obe casti maji jen jeden spolecny ramecek.
\tcbset{solutionlesson/.style={
  commonlesson,
  enhanced jigsaw,
  breakable,
  colbacklower=black!3,
  segmentation style={solid,draw=black!25,line width=0.35pt},
  fontlower=\small\color{notegray},
  before lower={\noindent\textbf{Řešení:}\enspace},
  leftlower=3mm,rightlower=3mm
}}
\newtcolorbox{prikladReseniBox}[1][]{solutionlesson,title={Příklad},#1}
\newtcolorbox{zadaniReseniBox}[1][]{solutionlesson,title={Zadání},#1}

% Pomucky pro autora pripravy: sedive poznamky; \reseni pro starsi pripravy
\newcommand{\poznamka}[1]{\par\begingroup\small\color{notegray}\itshape
  \textnormal{\textbf{Poznámka pro učitele.}}\enspace #1\par\endgroup}
\newcommand{\reseni}[1]{\par\begingroup\small\color{notegray}
  \textbf{Očekávané řešení:}\enspace #1\par\endgroup}
\newcommand{\blok}[1]{\needspace{5\baselineskip}\subsection*{#1}}
\newcommand{\vyrok}[1]{\uv{#1}}

\nastavPripravu{1. ročník}{Mocniny a odmocniny}{32}{Usměrňování zlomků}
\begin{document}
\poznamka{\textbf{45minutové jádro:} 4 min motivace; 8 min definice a Příklad 1;
8 min Samostatná práce 1 a)--d); 8 min strategie a Příklad 2;
8 min Samostatná práce 2 a)--d); 6 min proměnná a podmínky;
3 min závěr. Volitelné pokračování: třetí odmocnina, složitější proměnné a domácí práce.
\textbf{Dnes jmenovatel vždy obsahuje jen jeden odmocninový člen.}
Sdružené výrazy přijdou až v hodině 33.}

\blok{Proč se chceme zbavit odmocniny ve jmenovateli?}
Zlomky $1/\sqrt2$ a $\sqrt2/2$ mají stejnou hodnotu, ale ve druhém
je jmenovatel racionální. \textbf{Nesmíme změnit hodnotu zlomku}: násobíme
proto čitatele a jmenovatele stejným \emph{nenulovým} výrazem.
\[
\frac1{\sqrt2}\cdot\frac{\sqrt2}{\sqrt2}
=\frac{\sqrt2}{(\sqrt2)^2}=\boxed{\frac{\sqrt2}{2}}.
\]
\begin{definiceBox}[title={Definice: usměrňování zlomků}]
\textbf{Usměrňování zlomku} je úprava, při níž odstraníme odmocninu
ze jmenovatele, aniž změníme hodnotu výrazu.
Pro $a>0$ a libovolné reálné $k$ platí
\[
\boxed{\frac{k}{\sqrt a}=\frac{k\sqrt a}{a}}.
\]
Využíváme rozšíření zlomku výrazem $\sqrt a/\sqrt a=1$.
\end{definiceBox}
\poznamka{Nejde o nové pravidlo pro mocniny; jde o rozšíření zlomku.
Zeptej se, proč musí být $a>0$: pro $a=0$ by původní jmenovatel
byl nula; pro $a<0$ není druhá odmocnina reálná.}

\begin{prikladReseniBox}[breakable=false,title={Příklad 1: jednoduché odmocniny ve jmenovateli}]
Usměrni $\dfrac{5}{\sqrt3}$ a $\dfrac{3}{2\sqrt5}$.
\tcblower
\[
\frac5{\sqrt3}=\frac{5\sqrt3}{3},\qquad
\frac3{2\sqrt5}\cdot\frac{\sqrt5}{\sqrt5}
=\frac{3\sqrt5}{10}.
\]
Při druhém příkladu není nutné násobit čitatele a jmenovatele číslem $2$.
\end{prikladReseniBox}

\begin{zadaniReseniBox}[breakable=false,title={Samostatná práce 1: nejprve zautomatizuj postup}]
Usměrni, výsledek zkrať, lze-li:
\begin{enumerate}[label=\alph*)]
\item $\dfrac1{\sqrt5}$;
\item $\dfrac4{\sqrt7}$;
\item $\dfrac7{3\sqrt2}$;
\item $\dfrac6{5\sqrt3}$;
\item $\dfrac{2\sqrt3}{\sqrt5}$;
\item $\dfrac{3}{\sqrt{11}}$.
\end{enumerate}
\tcblower
\begin{enumerate}[label=\alph*)]
\item $\dfrac{\sqrt5}{5}$;
\item $\dfrac{4\sqrt7}{7}$;
\item $\dfrac{7\sqrt2}{6}$;
\item $\dfrac{6\sqrt3}{15}=\dfrac{2\sqrt3}{5}$;
\item $\dfrac{2\sqrt{15}}5$;
\item $\dfrac{3\sqrt{11}}{11}$.
\end{enumerate}
\end{zadaniReseniBox}

\blok{KISS: nejprve zjednoduš, teprve potom usměrni}
Po 27. hodině víme, že \textbf{správná strategie často zkrátí výpočet}.
Nejprve částečně odmocni nebo zjednoduš celý zlomek; až potom
rozhodni, zda je vůbec usměrňování potřeba.
\begin{prikladReseniBox}[breakable=false,title={Příklad 2: dva postupy, stejná hodnota}]
Usměrni $\dfrac4{\sqrt{12}}$ a uprav $\dfrac{\sqrt{45}}{\sqrt5}$.
\tcblower
\[
\frac4{\sqrt{12}}=\frac4{2\sqrt3}=\frac2{\sqrt3}
=\boxed{\frac{2\sqrt3}{3}}.
\]
\[
\frac{\sqrt{45}}{\sqrt5}=\sqrt{\frac{45}{5}}=\sqrt9=\boxed3.
\]
Ve druhém příkladu žádný iracionální jmenovatel \emph{po zjednodušení nezbyl}.
\end{prikladReseniBox}
\begin{zadaniReseniBox}[breakable=false,title={Samostatná práce 2: vyber nejkratší postup}]
\begin{enumerate}[label=\alph*)]
\item $\dfrac2{\sqrt8}$; \item $\dfrac9{\sqrt{27}}$;
\item $\dfrac{\sqrt{18}}{\sqrt2}$;
\item $\dfrac{\sqrt6}{\sqrt3}$;
\item $\dfrac{3\sqrt6}{2\sqrt3}$;
\item $\dfrac{12}{\sqrt{48}}$.
\end{enumerate}
\tcblower
\begin{enumerate}[label=\alph*)]
\item $\dfrac1{\sqrt2}=\dfrac{\sqrt2}{2}$;
\item $\dfrac3{\sqrt3}=\sqrt3$;
\item $3$; \item $\sqrt2$;
\item $\dfrac{3\sqrt2}2$;
\item $\dfrac3{\sqrt3}=\sqrt3$.
\end{enumerate}
\end{zadaniReseniBox}

\blok{Proměnné a podmínky nezmizí úpravou}
\begin{prikladReseniBox}[breakable=false,title={Příklad 3: usměrňování s proměnnou}]
Pro $x>0$ uprav $\dfrac3{\sqrt x}$ a $\dfrac{\sqrt{2x}}{\sqrt x}$.
\tcblower
\[
\frac3{\sqrt x}=\boxed{\frac{3\sqrt x}{x}},\qquad
\frac{\sqrt{2x}}{\sqrt x}=\sqrt{\frac{2x}{x}}=\boxed{\sqrt2}.
\]
Původní podmínka $x>0$ \textbf{platí i po zjednodušení}.
\end{prikladReseniBox}
\begin{zadaniReseniBox}[breakable=false,title={Samostatná práce 3: určuj i definiční obor}]
\begin{enumerate}[label=\alph*)]
\item $\dfrac{2}{\sqrt a}$;
\item $\dfrac{\sqrt b}{3\sqrt b}$;
\item $\dfrac5{\sqrt{3x}}$;
\item $\dfrac{\sqrt{12t}}{\sqrt{3t}}$.
\end{enumerate}
\tcblower
\begin{enumerate}[label=\alph*)]
\item $2\sqrt a/a$, $a>0$;
\item $1/3$, $b>0$ (původní jmenovatel nesmí být nulový);
\item $5\sqrt{3x}/(3x)$, $x>0$;
\item $2$, $t>0$.
\end{enumerate}
\end{zadaniReseniBox}

\blok{Rozšíření: usměrnění třetí odmocniny}
\poznamka{Následující pravidlo je mimo povinné jádro; nevytváří
předpoklad pro hodinu 33. Už známe třetí odmocninu z hodiny 29.}
\begin{prikladReseniBox}[breakable=false,title={Rozšiřující příklad 1: doplň do třetí mocniny}]
Usměrni $\dfrac1{\sqrt[3]4}$ a $\dfrac2{\sqrt[3]9}$.
\tcblower
Aby se $\sqrt[3]4=\sqrt[3]{2^2}$ změnila v celé číslo,
doplníme ještě jednu dvojku:
\[
\frac1{\sqrt[3]4}\cdot\frac{\sqrt[3]2}{\sqrt[3]2}
=\frac{\sqrt[3]2}{2}.
\]
U $\sqrt[3]9=\sqrt[3]{3^2}$ doplníme trojku:
$\dfrac2{\sqrt[3]9}=\dfrac{2\sqrt[3]3}{3}$.
\end{prikladReseniBox}

\begin{zadaniReseniBox}[breakable=false,title={Rozšiřující příklad 2: rozhodni, co násobit}]
\begin{enumerate}[label=\alph*)]
\item $\dfrac3{\sqrt[3]2}$;
\item $\dfrac1{\sqrt[3]{25}}$;
\item Proč postup $\dfrac1{\sqrt2+1}\cdot\dfrac{\sqrt2}{\sqrt2}$
\emph{neodstraní} odmocninu ze jmenovatele?
\end{enumerate}
\tcblower
\begin{enumerate}[label=\alph*)]
\item $\dfrac{3\sqrt[3]4}{2}$;
\item $\dfrac{\sqrt[3]5}{5}$;
\item Vznikne jmenovatel $2+\sqrt2$, v němž odmocnina zůstane.
Na dvojčleny použijeme v \textbf{hodině 33} jiný postup.
\end{enumerate}
\end{zadaniReseniBox}

\blok{Domácí procvičování}
\begin{zadaniReseniBox}[breakable=false,title={Domácí úkol 1: usměrni a kontroluj podmínky}]
\begin{enumerate}[label=\alph*)]
\item $\dfrac7{\sqrt{14}}$; \item $\dfrac5{2\sqrt3}$;
\item $\dfrac8{\sqrt{32}}$; \item $\dfrac{\sqrt{75}}{\sqrt3}$;
\item $\dfrac{\sqrt{50}}{\sqrt8}$;
\item $\dfrac{4}{\sqrt y}$ pro proměnnou $y$ (uveď podmínku).
\end{enumerate}
\tcblower
\begin{enumerate}[label=\alph*)]
\item $\sqrt{14}/2$;
\item $5\sqrt3/6$;
\item $\sqrt2$;
\item $5$;
\item $5/2$;
\item $4\sqrt y/y$, $y>0$.
\end{enumerate}
\end{zadaniReseniBox}
\poznamka{Závěrečná otázka: kdy lze rozšířit zlomek výrazem
$\sqrt a/\sqrt a$? Jen pro $a>0$, aby šlo o číslo 1.
Příště využijeme rozdíl čtverců při usměrňování jmenovatele s \emph{dvěma} členy.}
\end{document}
